\section{Relato de los hechos}\label{sec:relato}
Esta sección reconstruye en orden cronológico lo informado por organismos públicos, investigadores y prensa, con cada dato atribuido a su fuente. Las horas son locales (UTC$-$3) y aproximadas. Cuando las fuentes difieren, se presentan todas las versiones, sin resolverlas.

\subsection*{Antecedentes}
\paragraph{La cuenca y el abanico.} Según el análisis del Centro de Ciencia del Clima y la Resiliencia (CR2), Quebrada Honda drena una cuenca pequeña de la precordillera, de unos \SI{10,6}{km^2} y entre \SI{929}{m} y \SI{2355}{m} de altitud, que desemboca en el límite urbano de Santiago \citep{cr2_2026}. Según La Tercera, al pie se extiende el barrio de San Carlos de Apoquindo, donde se concentran colegios, clínicas y universidades \citep{latercera2026zonacero}. El decano de Arquitectura de la Universidad del Desarrollo, citado por El Desconcierto, explicó que varias quebradas de la ladera confluyen en un solo cauce y que «la quebrada se entuba en Bulnes Correa y el problema con los colectores cerrados es que tienen capacidad limitada» \citep{eldesconcierto2026}. Según dijo a The Clinic, ese cauce corre por un parque junto a General Blanche y, ya entubado, sigue bajo la calle Quebrada Honda \citep{theclinic2026allard}.

\paragraph{El incendio.} Según el CR2, el 29 de diciembre de 2025 comenzó un incendio en la cuenca \citep{cr2_2026}. El municipio informó que partió en el sector La Plaza Sur, que afectó una superficie preliminar de \SI{1012}{ha} en la comuna y que Conaf lo dio por controlado la tarde del 1 de enero \citep{lascondes2026}. Las cifras de superficie varían: «cerca de 1.000 hectáreas» según el CR2, que agrega que el fuego comprometió prácticamente toda la cuenca \citep{cr2_2026}; «aproximadamente mil hectáreas» en cuatro días según T13 \citep{t13_2026mallas}; y «más de 800 hectáreas», con el fuego a \SI{200}{m} de zonas urbanas, según Interferencia \citep{interferencia2026}. Según el cálculo de este estudio (sección~\ref{sec:res-incendio}), \SI{8,2}{km^2} de los \SI{10,1}{km^2} de las cuencas aportantes muestran señal de quema, con severidad mayoritariamente baja o moderada-baja.

\paragraph{El evento del 31 de enero.} Según el reporte de Sernageomin, ese día se activaron las quebradas La Honda, Grande, Los Almendros y Apoquindo \citep{sernageomin2026reterm}. Según T13, cayeron \SI{75}{mm} en una hora y \SI{20}{min}, y el evento afectó ocho hectáreas urbanas, casi \SI{4}{km} de calles y 11 viviendas \citep{t13_2026mallas}. Interferencia, en cambio, menciona «apenas 20 minutos de lluvia intensa» \citep{interferencia2026}. El Mostrador cifra en 22 las viviendas dañadas, según antecedentes iniciales de Senapred \citep{elmostrador2026}. La alcaldesa de Las Condes dijo entonces, según La Tercera, que el colector de Francisco Bulnes Correa «funcionó muy bien», aunque hubo puntos afectados, como las calles Mirasol y Carlos Peña \citep{latercera2026enero}.

\paragraph{La advertencia de Sernageomin.} Según el propio reporte, tras un sobrevuelo del 6 de febrero el Servicio Nacional de Geología y Minería elaboró el 11 de febrero el Reporte Técnico de Remoción en Masa Metropolitana 02/2026 \citep{sernageomin2026reterm}. Según Chilevisión, el informe situó el área en una zona de «alto peligro» por flujos de detritos e inundaciones, identificó flujos en las quebradas La Honda, Grande y Apoquindo, y señaló «deficiencias en las obras de mitigación y adecuación» \citep{chilevision2026}. El reporte de Sernageomin concluyó que el evento de enero «no constituye una ocurrencia aislada, sino una reactivación de procesos en un área altamente propensa» a estos fenómenos \citep{sernageomin2026reterm}. Según Chilevisión, el reporte recomendó, entre otras medidas, planes de evacuación y la limpieza periódica de los cauces \citep{chilevision2026}.

\paragraph{Las obras de mitigación.} Según Interferencia, el 26 de marzo la alcaldesa pidió al delegado presidencial instalar mallas dinámicas; él reconoció que sería difícil tenerlas antes de mayo \citep{interferencia2026}. Según el mismo medio, el 7 de agosto se anunció un convenio para cinco barreras con \$754 millones del Ministerio de Obras Públicas (MOP) y un plazo de 90 días \citep{interferencia2026}. T13 informa otra cifra: un contrato por trato directo de \$652 millones, firmado en septiembre, también a 90 días \citep{t13_2026mallas}. Según Interferencia, durante la emergencia el delegado presidencial confirmó que aún no estaban listas \citep{interferencia2026}.

\subsection*{La lluvia}
Según el CR2, un río atmosférico de categoría 2 a 3 alcanzó la zona central el 6 y el 7 de octubre. Dejó \SI{72}{mm} en Quinta Normal y cerca de \SI{100}{mm} en la precordillera de Santiago \citep{cr2_2026}. El 8 persistía una baja segregada estacionaria a unos \SI{5}{km} de altura que, según el CR2, junto con la humedad del río atmosférico, generó núcleos convectivos con oscurecimiento repentino, rayos, granizo y lluvia intensa en periodos de \SIrange{30}{60}{min} \citep{cr2_2026}. El CR2 señala que no parece haber pluviómetros en la cuenca; en la estación Río Molina, a unos \SI{7}{km}, cayeron \SI{42,8}{mm} ese día, \SI{31,6}{mm} de ellos entre las 15:00 y las 19:00, con un máximo de \SI{10,9}{mm} en \SI{30}{min}. Según ese análisis, la isoterma de \SI{0}{\celsius} se situó a unos \SI{2800}{m}, de modo que toda la cuenca recibió lluvia y no nieve \citep{cr2_2026}. Según el CR2, los rayos llegaron al nororiente de Santiago y a Las Condes cerca de las 16:30, y hacia las 17:00 se localizaban sobre todo en el Cajón del Maipo \citep{cr2_2026}. En Tobalaba, la Dirección Meteorológica de Chile registró \SI{34,8}{mm} el día 8 y \SI{13,8}{mm} entre las 15:00 y las 21:00 \citep{dmc2026}.

Otras cifras circularon sin medición identificada. La alcaldesa habló de \SI{144}{mm} acumulados, \SI{20}{mm} de ellos ese día, con el suelo saturado \citep{interferencia2026}. Según la versión consultada de Wikipedia, el subsecretario del Interior mencionó unos \SI{15}{mm} en un intervalo muy breve \citep{wikipedia2026}. Búsquedas en prensa registradas el 9 de octubre devolvieron también cifras de \SI{16}{mm} a \SI{20}{mm} en alrededor de una hora, sin estación identificada. La cifra de \SI{27}{mm} en \SI{20}{min} no se encontró en la prensa consultada y no está verificada. Los \SI{25}{mm} que también circularon parecen corresponder a un pronóstico para el 11 al 13 de octubre difundido por Canal 9, no a la lluvia caída el día 8 (búsqueda registrada el 12 de octubre).

\subsection*{El aluvión}
Cerca de las 17:00, según información preliminar de Carabineros, se divisó un «gran torrente de agua por calle General Blanche», que esa información atribuía al rebalse del canal de la calle Francisco Bulnes Correa; el caudal bajó hacia el poniente \citep{theclinic2026frontal,eldesconcierto2026}. Las fuentes ubican el inicio en puntos distintos. La Tercera informó que «todo comenzó al activarse una quebrada a la altura de General Blanche con San Carlos de Apoquindo», con avance hacia el surponiente \citep{latercera2026}. Según La Tercera, Carabineros, con un dron, ubicó la «zona cero» en Mons.\ Álvaro del Portillo con Av.\ La Plaza \citep{latercera2026zonacero}. Según Chilevisión, las autoridades atribuyeron el evento a la activación de la «Quebrada Las Ondas» \citep{chilevision2026mapa}, mientras Chilevisión, a partir del informe de Sernageomin, y el CR2 hablan de la quebrada La Honda u Honda \citep{chilevision2026,cr2_2026}.

El vicerrector de Comunicaciones de la Universidad de los Andes relató a Cooperativa un desborde de la quebrada en el lado norte del campus, que colinda con General Blanche, «por donde cayó toda esta masa de barro» \citep{cooperativa2026}. Según El Desconcierto, el flujo arrastró autos y un bus del sistema RED, cuyos pasajeros debieron subir al techo \citep{eldesconcierto2026}. La Tercera y T13 informaron que el bus fue volcado \citep{latercera2026,t13_2026vehiculos}; Cadena 3, que un vehículo de transporte público logró detenerse \citep{cadena3_2026}. Según el municipio, el barro bajó por Camino El Alba, General Blanche y Quebrada Honda \citep{latercera2026,chilevision2026mapa}. Según El Dínamo, Senapred señaló como más afectadas General Blanche, Quebrada Honda y la Av.\ La Plaza Sur \citep{eldinamo2026}. Según T13, el agua se acumuló también en la Av.\ Padre Hurtado \citep{t13_2026vehiculos}.

En la Clínica UC Christus, Cadena 3 informó cerca de \SI{1}{m} de barro en el servicio de urgencias \citep{cadena3_2026}. Según El Desconcierto, se suspendió la atención en ese servicio y las consultas \citep{eldesconcierto2026}. Según BioBioChile, en la sede del Duoc UC estudiantes y trabajadores quedaron atrapados en un auditorio \citep{biobio2026duoc}. Según Cooperativa, la comunidad de la Universidad de los Andes permaneció en el campus por recomendación de Bomberos \citep{cooperativa2026}. Un video difundido en Facebook se titula «Aluvión ingresó hasta el aula Magna de la Universidad de Los Andes» \citep{quintaconectado2026}, y Meganoticias publicó a las 18:16 imágenes del campus inundado \citep{meganoticias2026uandes}. Interferencia informó además agua y barro en el Hospital de la FACh \citep{interferencia2026}.

Sobre el drenaje se informaron versiones distintas, todas preliminares. Según Chilevisión, el delegado presidencial presumió «el bloqueo, el taponeo del canal subterráneo» en Álvaro del Portillo con Av.\ La Plaza \citep{chilevision2026}. Según Interferencia, en otra declaración el delegado apuntó a la obstrucción de dos ductos en General Blanche \citep{interferencia2026}. Según El Desconcierto, el canal entubado que baja desde la Av.\ La Plaza llega a una cámara de decantación en Francisco Bulnes Correa y, «con ese sistema colapsado, la calle funcionó como el nuevo cauce» \citep{eldesconcierto2026}. Interferencia subrayó que era una hipótesis sin confirmación técnica \citep{interferencia2026}.

\subsection*{Consecuencias y respuesta}
\paragraph{La tarde y la noche del 8.} Según Interferencia, Senapred dio su primer aviso público a las 17:41, con un video de su director regional; no envió un mensaje de evacuación a los celulares \citep{interferencia2026}. La alerta roja para diez comunas rigió desde las 18:20 según Interferencia \citep{interferencia2026}; La Tercera la sitúa a las 18:30 \citep{latercera2026lluvias}. Según La Tercera, a las 18:15 el presidente de la República escribió que se enfrentaba «una emergencia de activación muy rápida» \citep{latercera2026lluvias}. Según Interferencia, esa noche el presidente instruyó decretar zona de catástrofe tras un comité regional de gestión del riesgo \citep{interferencia2026}; según la versión consultada de Wikipedia, la medida se formalizó el 9 de octubre \citep{wikipedia2026}. El delegado presidencial informó, según Interferencia y El Desconcierto, seis lesionados, entre ellos un niño de un año con una fractura, y ningún fallecido \citep{interferencia2026,eldesconcierto2026}. Según Chilevisión, el ministro del Interior anunció un polígono de exclusión de ingreso tras un reconocimiento aéreo \citep{chilevision2026mapa}. Pasadas las 23:00, la alcaldesa dijo a Cooperativa que «el cerro cedió de manera casi completa» \citep{cooperativa2026alcaldesa}.

\paragraph{El perímetro.} Según The Clinic, el primer perímetro municipal de seguridad iba de Camino El Alba a Quebrada Honda, con Vital Apoquindo y el límite cordillerano por el oriente, sujeto a ajustes \citep{theclinic2026frontal}. Según La Tercera, el polígono informado después tenía por límites Camino El Alba (norte), la Av.\ La Plaza (oriente), Quebrada Honda (sur) y la Av.\ Padre Hurtado (poniente), en el sector del Pueblito Los Dominicos \citep{latercera2026}.

\paragraph{Superficie y viviendas.} Las cifras difieren según el organismo y la fecha:
\begin{itemize}
\item Superficie: unas \SI{200}{ha} según el balance de la noche del 8 \citep{wikipedia2026}; \SI{253}{ha} y 3.981 roles de avalúo según la plataforma inmobiliaria Inciti \citep{chicureohoy2026}; \SI{288,5}{ha}, 2.266 predios y 368 tramos de calle en el visor oficial de IDE Chile \citep{bbnn2026visor}.
\item Municipalidad (9 de octubre): 2.015 viviendas afectadas, 55 inhabitables y 1.960 con daños menores, según CNN Chile y Chicureo Hoy \citep{cnnchile2026balance,chicureohoy2026}. Ese mismo día, el delegado presidencial habló de 1.960 viviendas afectadas, 55 de ellas con «problemas serios de habitabilidad», según Publimetro \citep{publimetro2026}.
\item Senapred: 75 viviendas con daño mayor y 2.286 con daño menor el 9 de octubre, según La Tribuna \citep{latribuna2026}; más tarde, según La Tercera, 56 destruidas, 40 con daño mayor, 1.960 con daño menor y 384 damnificados \citep{latercera2026senapred}.
\item Visor de Bienes Nacionales: 22 viviendas con daño mayor, 1.960 con daño menor y 51 equipamientos afectados \citep{bbnn2026visor}.
\end{itemize}
La municipalidad cifró en 1.200 las personas evacuadas o rescatadas, según CNN Chile \citep{cnnchile2026balance}.

\paragraph{La Clínica Universidad de los Andes.} Las versiones no coinciden. Meganoticias informó que el barro ingresó hasta el área de urgencias \citep{meganoticias2026urgencia}, y BioBioChile difundió imágenes de salas y pasillos anegados atribuidas a esta clínica y a la Clínica UC Christus \citep{biobio2026clinicas}. La directora de Comunicaciones de la clínica afirmó a Radio La Señal que «toda la afectación fue en el exterior» y que solo entró «un poco de agua en una sala de espera», de no más de un centímetro \citep{radiolasenal2026}. El balance municipal consignó «un ingreso menor de agua», con la urgencia disponible, según CNN Chile \citep{cnnchile2026balance}.

\paragraph{Las mallas.} Según T13, el 9 de octubre el ministro de Obras Públicas dijo que había habido «instalaciones de mallas dinámicas». La alcaldesa precisó luego que faltaba habilitar los accesos y que la instalación «partía mañana»; el ministro reconoció después que los trabajos estaban en una etapa previa \citep{t13_2026mallas}. Según Chicureo Hoy, el Ministerio Público abrió una investigación sobre eventuales responsabilidades penales \citep{chicureohoy2026}.

\subsection*{Preguntas que dejan los hechos}
No hay acuerdo sobre el punto por donde el flujo salió a la ciudad: Álvaro del Portillo con La Plaza, General Blanche con San Carlos de Apoquindo o el canal de Francisco Bulnes Correa. Tampoco está establecido si el barro provino de un solo cauce o de varias quebradas, como sugirió a CNN Chile un concejal de Las Condes \citep{cnnchile2026haddad}. No se localizaron en las fuentes consultadas mediciones de la lluvia dentro de la cuenca ni del volumen movilizado, que el delegado presidencial describió como «miles de metros cúbicos» \citep{cooperativa2026alcaldesa}. Las superficies y conteos publicados difieren entre sí y ninguno informa espesores ni marcas de nivel. Estas preguntas motivan las hipótesis de la sección~\ref{sec:hipotesis}.
